% Сборка: pdflatex -interaction=nonstopmode -halt-on-error shared_polynomial_edr.tex (дважды).
\documentclass[10pt,a4paper]{article}
\usepackage[T2A]{fontenc}
\usepackage[utf8]{inputenc}
\usepackage[english,russian]{babel}
\usepackage[margin=20mm]{geometry}
\usepackage{amsmath,amssymb,microtype}
\usepackage[hidelinks]{hyperref}
\usepackage{xurl}
\setlength{\parindent}{0pt}
\setlength{\parskip}{3pt}
\setlength{\emergencystretch}{2em}
\newcommand{\R}{\mathbb R}
\newcommand{\norm}[1]{\lVert#1\rVert}
\newcommand{\code}[1]{\nolinkurl{#1}}
\begin{document}
\begin{center}
{\Large\bfseries Общее полиномиальное приближение для EDR}\\[3pt]
{\small Экспериментальный метод: функционал, алгоритм и ограничения}
\end{center}

\paragraph{Постановка.}
Даны $x_i\in\R^d$, $y_i\in\R$, $i=1,\dots,n$, и задана размерность $m<d$.
Предполагается многоиндексная модель
\[
 y_i=f(B_*^\top x_i)+\varepsilon_i,\qquad
 B_*\in\R^{d\times m},\quad B_*^\top B_*=I_m.
\]
Цель --- оценить подпространство $\operatorname{span}(B_*)$; функция $f$ и
истинный базис алгоритму неизвестны. Для ортонормированного кандидата
$B\in\R^{d\times m}$ полагаем $z_i=B^\top x_i$ и приближаем отклик одним
общим полиномом степени не выше $q=3$:
\[
 p_c(z)=\sum_{|\alpha|\le 3}c_\alpha h_\alpha(z),\qquad
 h_\alpha(z)=\prod_{k=1}^m
 \frac{\operatorname{He}_{\alpha_k}(z_k)}{\sqrt{\alpha_k!}}.
\]
Здесь $\operatorname{He}_j$ --- вероятностный полином Эрмита; входят \emph{все}
смешанные члены. Число коэффициентов $p=\binom{m+3}{3}$: при $m=3$ их 20,
при $m=5$ --- 56. Это общая модель для всей выборки, а не набор независимых
локальных линейных моделей.

\paragraph{Минимизируемый функционал.}
Пусть $D_B\in\R^{n\times p}$, $(D_B)_{i\alpha}=h_\alpha(B^\top x_i)$, и
$y=(y_1,\ldots,y_n)^\top$. Оптимизируется профиль
\begin{equation}\label{eq:profile}
 \boxed{\displaystyle
 L(B)=\frac1n\min_{c\in\R^p}\norm{D_Bc-y}^2,
 \qquad B^\top B=I_m.}
\end{equation}
Для каждого пробного $B$ коэффициенты $\hat c(B)$ находятся методом
наименьших квадратов через экономное QR-разложение $D_B=QR$, где
$Q^\top Q=I_p$: $R\hat c=Q^\top y$. При полном столбцовом ранге
$L(B)=n^{-1}\norm{(I_n-QQ^\top)y}^2$. Постоянного ridge-штрафа на
коэффициенты полинома нет. Полный набор полиномов делает значение профиля
инвариантным к ортогональному повороту базиса $B\mapsto BO$.

\paragraph{Градиент профиля.}
Определим $r_i=p_{\hat c}(B^\top x_i)-y_i$ и
$g_i=\nabla_zp_{\hat c}(B^\top x_i)\in\R^m$. При полном ранге $D_B$,
по теореме об огибающей, горизонтальный градиент равен
\begin{equation}\label{eq:gradient}
 \operatorname{grad}L(B)=\frac{2}{n}(I_d-BB^\top)
 \sum_{i=1}^n r_i x_i g_i^\top.
\end{equation}
Проекция $I_d-BB^\top$ убирает повороты внутри самого подпространства.
Уменьшается ошибка отклика $L(B)$; истинная EDR при обучении не вычисляется.

\newpage
\paragraph{Анизотропный шаг Schur--Gauss--Newton.}
Начальный базис $B_0$ получается полным ADP с метрикой \code{two_sided}.
На каждой итерации выбирается ортонормированное дополнение
$V\in\R^{d\times(d-m)}$, и горизонтальное приращение записывается как
$\Delta=VK$, $K\in\R^{(d-m)\times m}$. Пусть $t=m(d-m)$,
а $A\in\R^{n\times t}$ --- матрица для построчной развёртки $K$, для которой
\[
 (A\operatorname{vec}K)_i=g_i^\top K^\top V^\top x_i.
\]
После исключения линейных приращений коэффициентов полинома строятся
\begin{equation}\label{eq:schur}
 H=\frac1nA^\top(I_n-QQ^\top)A,\qquad
 b=\frac1nA^\top r,\qquad
 (H+\rho I_t)\operatorname{vec}K=-b.
\end{equation}
В коде используется $\rho=10^{-4}\max\{\operatorname{tr}(H)/t,10^{-12}\}$.
Демпфирование относится к \emph{шагу}, а не является штрафом в
\eqref{eq:profile}. Приращение ограничивается условием
$\norm{VK}_2\le0.5$; для $B_{\rm trial}=\operatorname{polar}(B+\eta VK)$
поиск по длине шага требует уменьшения исходного профиля. Число итераций
ограничено 60, порог нормы градиента равен $10^{-7}$.

Матрица $H$ --- Schur-приближение совместного шага Гаусса--Ньютона после
исключения приращений $c$. При ненулевом остатке оно не является точным
якобианом профилированного остатка. Проверяются численный ранг матрицы $D_B$,
остаток нормальных уравнений и ранги градиентного Gram и $H$; это численные
проверки, а не сертификат статистической идентифицируемости.

\paragraph{Связь с истинной EDR.}
Пусть $X\sim N(0,I_d)$, $E[Y\mid X]=f(B_*^\top X)$ и $f_{\le3}$ ---
проекция этой функции на Gaussian-chaos порядков не выше трёх. Определим
\[
 G_3=E[\nabla_x f_{\le3}(X)\nabla_x f_{\le3}(X)^\top],\qquad
 \operatorname{EDR}^2(B,B_*)=
 \frac{\norm{BB^\top-B_*B_*^\top}_F^2}{2m}.
\]
Если $\operatorname{rank}(G_3)=m$ и $\operatorname{range}(G_3)=
\operatorname{span}(B_*)$, то истинное подпространство --- единственный
глобальный минимум \emph{популяционного} риска $R_3(B)$. При этом
\begin{equation}\label{eq:bound}
 R_3(B)-R_3(B_*)\ge
 \frac{m\lambda_{\min}^{+}(G_3)}{3}\operatorname{EDR}^2(B,B_*).
\end{equation}
Слабое минимальное собственное значение ослабляет оценку. Формула не делает
выборочный профиль равным EDR и не гарантирует глобальной сходимости.

\paragraph{Ограничение и проверка.}
Если $f(X)=X_1+\delta h_4(X_2)$ при стандартном Gaussian $X$, кубическая
проекция теряет второе направление: её популяционный риск одинаков для
$\operatorname{span}(e_1,e_2)$ и $\operatorname{span}(e_1,e_3)$. Поэтому
метод нельзя считать универсальным. На независимой тестовой серии
$d=25$, $n=5000$, $m\in\{3,5\}$, три функции и 60 парных наборов средняя
EDR снизилась с $0.309$ у \code{two_sided} до $0.091$; при additive и $m=5$
восстановлений с EDR $\le0.1$ осталось $0/10$. На отдельном примере $h_4$
кубическое уточнение ухудшило EDR. Проверенный метод остаётся
экспериментальным, без изменения штатных настроек ADP.

\smallskip
\textbf{Реализация и доказательства:} \code{experiments/shared_polynomial.py};
\code{verification/shared-polynomial-20261028/THEORY.md},
\code{R1_AUDIT.md} и \code{REPORT.md} в том же каталоге.
\end{document}
